\section{Limites, menaces à la validité, suite}\label{sec:limites}

\subsection{Ce qui n'est pas établi}
\begin{enumerate}
\item \textbf{Attribution causale du gain.} Le passage de 0,5072 à 0,5700 mêle quatre changements : la représentation (unités, tokens, parcours d'ordres), l'architecture, l'ensemble de 4 modèles et le refit. Aucune soumission n'isole l'un d'eux. Le fait que \code{control\_gru}, qui reçoit les nouvelles entrées, dépasse nettement le GRU de R1 en validation suggère que la représentation pèse lourd. Mais la validation est biaisée.
\item \textbf{Validation non représentative.} Faute de dates, aucune partition ne tient des jours entiers à l'écart. Toutes les mesures internes sont optimistes. Les intervalles bootstrap sous-estiment la variance d'un facteur $\sqrt{\mathrm{deff}}$ inconnu.
\item \textbf{Leaderboard public.} Il est mesuré sur une partie du test, inconnue ici. Le consulter trop souvent ferait sur-apprendre ce sous-ensemble : \code{LEADERBOARD.md} limite et documente les soumissions.
\item \textbf{Hypothèse du prix.} « Prix plus élevés dans la période test » relie trois observations (profondeur, odd lots, lots ronds), mais le niveau de prix n'est pas observé.
\item \textbf{Transductivité.} Le sens du stress, l'équilibrage et le lissage par voisins utilisent les entrées test. Il faut vérifier que le règlement l'autorise. Dans le cas contraire, la soumission de référence est 0,5700.
\end{enumerate}

\subsection{Prochaines expériences}
\begin{center}\small
\begin{tabularx}{\linewidth}{>{\raggedright\arraybackslash}p{3.4cm}>{\raggedright\arraybackslash}X>{\raggedright\arraybackslash}p{3.8cm}}
\toprule
Expérience & Hypothèse & Abandon si\\
\midrule
V4 : pseudo-étiquetage & apprendre sur des fenêtres test étiquetées par V3 B (accord lissé/brut, quota par titre) adapte le modèle à la période & C\_k40 $\le$ D\_k40 au LB\\
V4 : 90 époques & l'entraînement n'était pas saturé à 60 & gain inférieur à l'écart entre graines\\
V4 : $k$ au test & $k=40$ n'est pas optimal & variantes à $\pm0{,}3$ pt de B\\
Token $D^{\text{pré}}_t$ (§\ref{sec:dpost}) & distinguer « où était l'ordre » de « ce qu'il a fait au carnet » & stress équilibré inchangé\\
SignatureLite & 3 convolutions + pooling (≈ 100\,k paramètres) suffisent, sans attention & écart > 0,5 pt vs \code{v3\_long}\\
Regroupement explicite en jours & voter par paquet de 20 fenêtres plutôt que par $k$ voisins & pas de gain sur B\\
\bottomrule
\end{tabularx}
\end{center}

\subsection{Conclusion}
Le gain le plus robuste de ce travail ne vient pas d'une couche de réseau. Il vient de trois décisions. D'abord, \emph{regarder} la donnée dans ses propres unités, ce qui a fait apparaître le régime de tick comme signature stable. Ensuite, \emph{refuser} de croire la validation aléatoire. Enfin, \emph{diagnostiquer le test sans labels} avant de le soumettre. L'équilibrage de Sinkhorn n'est qu'une conséquence de ce diagnostic : une projection entropique de deux lignes, dont la justification tient en une formule de Bayes.
